\chapter{Objetivos, Alcance y Antecedentes}
\label{cap:objetivos_alcance_antecedentes}

% ==============================================================================
% 2.1. OBJETIVOS DEL PROYECTO
% ==============================================================================
\section{Objetivos del proyecto}
\label{sec:objetivos_proyecto}

\subsection{Objetivo General}
El objetivo principal consiste en evaluar la viabilidad técnica, la robustez matemática y la competitividad financiera de los algoritmos cuánticos híbridos variacionales frente a los métodos clásicos de optimización en la resolución del problema de selección de carteras sujeto a restricciones discretas duras de cardinalidad exacta ($K$).

Con este propósito, la investigación busca determinar en qué medida el confinamiento geométrico del espacio de búsqueda y la estabilización del bucle variacional permiten superar las ineficiencias del paradigma estándar de penalizaciones cuadráticas, estableciendo de forma cuantitativa las capacidades y fronteras de escalabilidad de la computación cuántica en su estado actual frente a los estándares de la industria financiera.

\subsection{Objetivos Específicos}
\label{sec:objetivos_especificos}
Para dar cumplimiento estructurado y medible al objetivo general, se establecen los siguientes objetivos específicos:

\begin{itemize}

    \item \textbf{OE1. Verificar la equivalencia matemática del mapeo
    Markowitz $\rightarrow$ QUBO $\rightarrow$ Ising.}
    Comprobar que la energía del Hamiltoniano de Ising construido reproduce
    exactamente la función objetivo financiera original (media-varianza con
    penalización de cardinalidad) para el conjunto completo de
    configuraciones binarias en instancias de tamaño reducido ($N \le 14$),
    empleando un solucionador exacto por fuerza bruta como referencia de
    verificación. \textit{Criterio de éxito:} discrepancia numérica nula
    (dentro de la tolerancia de precisión de punto flotante) entre la energía
    de Ising y el valor QUBO en el 100\% de las configuraciones evaluadas.

    \item \textbf{OE2. Cuantificar la mejora en la factibilidad estructural
    de las soluciones cuánticas.}
    Determinar si la sustitución del mezclador transversal ($X$-Mixer) y la
    penalización cuadrática por un operador de intercambio $XY$ inicializado
    en un estado de Dicke incrementa de forma sustancial la proporción de
    muestras que satisfacen la restricción $\sum_i x_i = K$, frente al QAOA
    estándar, para un conjunto de instancias con $N \in [6, 20]$ y
    profundidades $p$ acotadas. \textit{Criterio de éxito:} diferencia
    estadísticamente relevante y creciente con $N$ entre la tasa de
    factibilidad de ambos enfoques, verificable mediante el muestreo directo
    del vector de estado simulado.

    \item \textbf{OE3. Evaluar el efecto estabilizador de la inicialización
    TQA y la regularización Ridge $L_2$ sobre la convergencia clásica.}
    Medir si la combinación de una inicialización basada en \textit{Trotterized
    Quantum Annealing} y una penalización continua $L_2$ sobre la
    trayectoria de los parámetros $(\gamma, \beta)$ reduce la variabilidad
    del resultado final del optimizador COBYLA y/o el número de evaluaciones
    de la función de coste necesarias para alcanzar un óptimo local estable,
    en comparación con una inicialización aleatoria sin regularizar, sobre
    una instancia de referencia fija. \textit{Criterio de éxito:} reducción
    observable de la desviación estándar del \textit{Optimization GAP} entre
    semillas aleatorias al introducir la regularización.

    \item \textbf{OE4. Contrastar el rendimiento del XY-QAOA regularizado
    frente a solucionadores de referencia clásicos.}
    Comparar el \textit{Optimization GAP} respecto al óptimo exacto
    (Gurobi), el Ratio de Sharpe fuera de muestra y el coste temporal de
    ejecución del XY-QAOA regularizado frente a Gurobi Optimizer y Simulated
    Annealing, utilizando series históricas reales de mercado, para
    determinar el rango de tamaños de universo $N$ en el que la propuesta
    cuántica resulta competitiva bajo simulación clásica exacta.
    \textit{Criterio de éxito:} identificación explícita del umbral de $N$ a
    partir del cual el coste de simulación deja de ser práctico, y
    caracterización cuantitativa de la relación GAP–robustez financiera en
    el rango accesible.

\end{itemize}

\subsection{Hipótesis de Trabajo}
\label{sec:hipotesis}

A partir del problema planteado, del análisis del estado del arte y de las limitaciones
documentadas en las formulaciones estándar de QAOA para problemas combinatorios con
restricciones, se establecen las siguientes hipótesis de trabajo, que guían el diseño
experimental del proyecto y que se contrastan explícitamente en la
Sección~\ref{sec:cumplimiento_objetivos} del Capítulo~\ref{chap:conclusiones}:

\begin{itemize}
    \item \textbf{H1.} La incorporación explícita de la restricción de cardinalidad en la
    propia dinámica del circuito (inicialización en estado de Dicke + mezclador $XY$)
    produce una tasa de factibilidad estructural superior a la del QAOA estándar con
    penalización cuadrática, y esta ventaja se acentúa a medida que crece $N$.

    \item \textbf{H2.} Una inicialización de los parámetros variacionales basada en
    \textit{Trotterized Quantum Annealing} (TQA) proporciona un punto de partida más
    informado que una inicialización aleatoria, reduciendo la variabilidad entre semillas
    del resultado final del optimizador clásico.

    \item \textbf{H3.} La incorporación de una penalización continua Ridge $L_2$ sobre la
    trayectoria de los parámetros aporta, de forma aislada, una mejora adicional y
    consistente sobre la estabilidad o calidad de la solución respecto a usar únicamente
    inicialización TQA.

    \item \textbf{H4.} Un XY-QAOA regularizado, evaluado bajo simulación clásica exacta y
    un único reinicio por semilla, alcanza una calidad de solución (GAP de optimización) y
    una robustez financiera fuera de muestra (Sharpe OOS) competitivas frente a
    solucionadores clásicos de referencia --el óptimo exacto (Gurobi) y una metaheurística
    estocástica simple (Simulated Annealing).

    \item \textbf{H5.} Para una profundidad de circuito $p$ fija, el ansatz XY-QAOA agota
    su capacidad de representación al crecer $N$, lo que se traduce en un repunte del GAP
    en las instancias de mayor tamaño estudiadas.
\end{itemize}

Estas hipótesis no se plantean como resultados asumidos de antemano, sino como
proposiciones que se contrastan mediante el diseño experimental descrito en el
Capítulo~\ref{ch:desarrollo} y los resultados del Capítulo~\ref{ch:resultados}. El grado
de confirmación de cada una se resume explícitamente en la
Sección~\ref{sec:cumplimiento_objetivos}.

% ==============================================================================
% 2.2. ALCANCE Y DELIMITACIÓN DEL ESTUDIO
% ==============================================================================
\section{Alcance}

El alcance del proyecto constituye un marco fundamental para comprender su propósito, 
los requerimientos técnicos identificados y los límites operacionales del estudio, 
estableciendo unas expectativas claras sobre los objetivos abordados.

\subsection{Dentro del Alcance}
\label{subsec:dentro_alcance}

El alcance propuesto incluye los siguientes puntos:
\begin{itemize}
    \item Revisión de la literatura y del estado del arte en algoritmos cuánticos 
          híbridos variacionales aplicados a problemas de optimización combinatoria 
          en finanzas.
    \item Recopilación, limpieza y preprocesamiento de series temporales de datos 
          históricos reales procedentes de mercados organizados (activos del S\&P 500 
          e IBEX 35).
    \item Formulación matemática del problema de selección de carteras bajo 
          restricción dura de cardinalidad exacta ($K$) y su mapeo a estructuras QUBO 
          y Hamiltonianos de espín de Ising.
    \item Desarrollo e implementación de la simulación cuántica variacional 
          utilizando circuitos con preservación de simetría estructural (estados de 
          Dicke y operador mezclador XY) sobre el framework Qrisp.
    \item Refinamiento del modelo e hiperparametrización del bucle híbrido variacional, 
          integrando inicialización adiabática troterizada (TQA) y regularización continua 
          Ridge $L_2$.
    \item Evaluación experimental y benchmarking sistemático de las soluciones obtenidas 
          frente a solucionadores clásicos deterministas de grado industrial (Gurobi Optimizer) 
          y metaheurísticas clásicas (Simulated Annealing).
\end{itemize}

\subsection{Fuera del Alcance}
\label{subsec:fuera_alcance}

En este sentido, se define formalmente lo que queda fuera del alcance de este trabajo:
\begin{itemize}
    \item La ejecución directa en hardware u ordenadores cuánticos reales (QPUs), 
          limitando la experimentación a la simulación matemática clásica exacta 
          mediante vectores de estado.
    \item La simulación de universos de activos con cardinalidad y dimensión muy 
          elevadas ($N > 20$), debido a la barrera exponencial en los requisitos de 
          memoria RAM y tiempos de cómputo de la simulación clásica.
    \item El desarrollo de modelos dinámicos multiperiodo con rebalanceo continuo y 
          costes de transacción por deslizamiento (el estudio se enfoca en la asignación 
          estática en un único periodo).
    \item La integración y procesamiento de flujos de datos de cotización en tiempo real 
          o estrategias de negociación de alta frecuencia (HFT).
    \item La implementación de esquemas activos de corrección de errores cuánticos (QEC).
\end{itemize}

% ==============================================================================
% 2.3. ANTECEDENTES Y ESTADO DEL ARTE
% ==============================================================================
\section{Antecedentes y estado del arte}
\label{sec:antecedentes_estado_arte}

\subsection{De Markowitz al problema combinatorio}
\label{subsec:markowitz_combinatorio}

Como se ha introducido en el Capítulo \ref{ch:introduccion}, la incorporación
de la restricción de cardinalidad exacta $K$ traslada el modelo de Markowitz
desde la clase de complejidad P hacia la clase NP-hard, con un espacio de
búsqueda discreto gobernado por el coeficiente binomial $\binom{N}{K}$ frente
al espacio completo $2^N$, tal como se cuantificó en la Tabla
\ref{tab:dimension_espacio}. Ante este reto computacional, la industria
financiera recurre tradicionalmente a dos familias de solucionadores
clásicos, cuyas fortalezas y limitaciones se resumen a continuación.

\subsection{Solucionadores clásicos de referencia}
\label{subsec:solucionadores_clasicos}

En la práctica cuantitativa actual, la resolución de modelos MIQP de selección de carteras se fundamenta en dos aproximaciones complementarias:

\begin{itemize}
    \item \textbf{Solucionadores Deterministas Exactos (Gurobi Optimizer)}
    \cite{Gurobi2025}: certifican la optimalidad global de la solución
    mediante planos de corte y bifurcación ramificada, al coste de un
    escalado temporal exponencial en el peor caso.

    \item \textbf{Metaheurísticas Estocásticas (Simulated Annealing)}: algoritmos
    de búsqueda local inspirados en el recocido de materiales, rápidos y
    escalables, pero sin certificado de convergencia global ni garantía de
    reproducibilidad exacta entre semillas pseudoaleatorias.
\end{itemize}

\subsection{Computación cuántica variacional}
\label{subsec:qaoa_evolucion}

La incapacidad de los métodos clásicos deterministas para mitigar el crecimiento combinatorio en universos densos impulsó el interés por la computación cuántica. Históricamente, las primeras aproximaciones prácticas se desarrollaron sobre plataformas analógicas de recocido cuántico (\textit{Quantum Annealing}), diseñadas para minimizar funciones de energía mediante tunelamiento cuántico en redes de espines físicos. No obstante, las restricciones de conectividad en los grafos de acoplamiento físico y la susceptibilidad al ruido ambiental motivaron la transición hacia el paradigma cuántico digital basado en compuertas lógicas parametrizadas.

En este marco digital de la era NISQ \cite{Preskill2018}, el Algoritmo de Optimización Cuántica Aproximada (QAOA), introducido por Farhi et al. \cite{Farhi2014}, se posicionó como el referente fundamental para optimización combinatoria. QAOA opera alternando la aplicación de operadores unitarios parametrizados por ángulos $(\gamma, \beta)$ correspondientes a un Hamiltoniano de coste (que codifica la función objetivo) y un Hamiltoniano mezclador transversal ($X$-Mixer). En la formulación estándar, la observancia de restricciones duras como la cardinalidad $K$ se gestiona mediante términos de penalización cuadrática integrados en la función de coste, mapeando el problema a un formato de Optimización Binaria Cuadrática No Restringida (QUBO).

El flujo general del algoritmo híbrido variacional se ilustra en la Figura \ref{fig:flujo_qaoa}.

\begin{figure}[htbp]
    \centering
    \vspace*{0.5cm}
    $\text{Inicialización de parámetros } (\gamma_0, \beta_0) \longrightarrow \framebox{Ejecución Circuito en QPU} \longrightarrow \framebox{Medición del Estado} \longrightarrow \framebox{Optimización Clásica (CPU)} \longrightarrow \text{Actualización de } (\gamma_{k+1}, \beta_{k+1})$
    \vspace*{0.5cm}
    \caption{Diagrama de flujo del bucle de optimización híbrido clásico-cuántico variacional (QAOA).}
    \label{fig:flujo_qaoa}
\end{figure}


\subsection{Limitaciones documentadas del QAOA estándar}
\label{subsec:limitaciones_qaoa}

El QAOA estándar con penalizaciones
presenta dos limitaciones estructurales ampliamente documentadas en la
literatura especializada cuando se aplica a problemas fuertemente
restringidos: (i) la deformación del paisaje energético inducida por
multiplicadores de penalización de gran magnitud, y (ii) el fenómeno de
mesetas estériles o \textit{Barren Plateaus}, por el cual la varianza del
gradiente de la función de coste se desvanece con el número de cúbits del
registro \cite{McClean2018}. El mecanismo formal de ambas patologías, así
como la arquitectura propuesta en este trabajo para superarlas, se
desarrolla en detalle en la Sección \ref{subsec:barren_plateaus} del
Capítulo \ref{ch:marco_teorico}.

\subsection{Trabajos relacionados: mezcladores XY, inicialización TQA y QAOA en carteras}
\label{subsec:trabajos_relacionados}

La restricción del espacio de búsqueda mediante mezcladores que preservan el peso de
Hamming fue formalizada dentro del marco general del \textit{Quantum Alternating
Operator Ansatz}\cite{Hadfield2019}, y su aplicación específica a mezcladores $XY$ inicializados en estados
de Dicke ha sido validada tanto en simulación como en hardware real de trampa de iones
por He et al. \cite{He2023}, quienes muestran que la alineación entre el estado inicial y
el mezclador mejora sistemáticamente el rendimiento de QAOA en problemas con
restricciones. Niroula et al. \cite{Niroula2022} aplican un enfoque análogo a un problema
de resumen extractivo, confirmando la generalidad de esta familia de mezcladores más allá
de las finanzas. Brandhofer et al. \cite{Brandhofer2022} realizan un \textit{benchmarking}
sistemático de QAOA aplicado a optimización de carteras, comparando distintas
formulaciones, tamaños de universo y estrategias de penalización frente a solucionadores
clásicos, en línea metodológica con el enfoque comparativo adoptado en este trabajo.

Frente al trabajo de Mancilla et al. \cite{Mancilla2026} --la propuesta más cercana en la
literatura--, este TFM aporta una caracterización experimental más sistemática del punto
de ruptura entre viabilidad estructural y competitividad en calidad de solución, en lugar
de limitarse a un único escenario de mercado favorable.

El trabajo más directamente comparable a esta propuesta es el de Mancilla et al.
\cite{Mancilla2026}, quienes combinan igualmente inicialización en estado de Dicke,
mezclador $XY$ e inicialización troterizada de tipo adiabático para la selección de
carteras bajo restricción de cardinalidad, evaluando su enfoque mediante \textit{backtesting}
sobre datos reales de mercado. A diferencia de dicho trabajo, que reporta una ventaja de
Sharpe Ratio del enfoque cuántico frente a Simulated Annealing en un backtest de
\textit{walk-forward} sobre 10 activos con una única ventana temporal, este TFM evalúa el
comportamiento del algoritmo bajo un barrido sistemático de tamaño de universo
($N=6$ a $N=20$) y de profundidad de circuito, y bajo el régimen de un único reinicio por
semilla —más representativo del coste de ejecución en hardware NISQ real que el mejor de
varios reinicios—, lo que permite caracterizar de forma más granular el punto exacto en el
que la ventaja de factibilidad estructural del enfoque cuántico se traduce, o no, en una
ventaja de calidad de solución frente a una metaheurística clásica simple.

Por último, la estrategia de inicialización basada en \textit{Trotterized Quantum
Annealing} adoptada en la Sección~\ref{sec:estabilizacion_optimizacion} del
Capítulo~\ref{ch:marco_teorico} tiene su origen en el trabajo de Sack y Serbyn
\cite{SackSerbyn2021}, quienes demuestran que esta inicialización permite evitar mínimos
falsos en el paisaje de optimización de QAOA para un amplio rango de pasos temporales de
Trotter, con un rendimiento comparable al mejor resultado obtenido tras un número
exponencialmente creciente de inicializaciones aleatorias.

% ==============================================================================
% 2.4. SÍNTESIS COMPARATIVA Y POSICIONAMIENTO DEL TRABAJO
% ==============================================================================
\section{Posicionamiento del trabajo}
\label{sec:sintesis_posicionamiento}

Frente a las limitaciones documentadas del QAOA estándar y la ineficiencia de evaluar el espacio completo de Hilbert $2^N$, este trabajo adopta un enfoque de diseño algorítmico avanzado. En lugar de sancionar las soluciones no factibles mediante penalizaciones algebraicas externas, se restringe la dinámica física del circuito cuántico al subespacio geométrico estricto de las carteras válidas $\binom{N}{K}$, complementando dicha estructura con un esquema bicapa de inicialización adiabática guiada (TQA) y regularización Ridge $L_2$ en el optimizador clásico.

La Tabla~\ref{tab:estado_arte} sintetiza la comparativa taxonómica entre los diferentes paradigmas de optimización clásica y cuántica analizados en este estado del arte, posicionando conceptualmente la arquitectura propuesta dentro del panorama tecnológico contemporáneo.

\begin{table}[htbp]
    \centering
    \small
    \caption{Comparativa del estado del arte en algoritmos de optimización combinatoria clásica y cuántica.}
    \label{tab:estado_arte}
    \renewcommand{\arraystretch}{1.3}
    \begin{tabular}{p{2.6cm}p{2.4cm}p{2.8cm}p{2.2cm}p{2.7cm}}
        \toprule
        \textbf{Algoritmo / Enfoque} & \textbf{Paradigma de Computación} & \textbf{Manejo de Restricciones ($K$)} & \textbf{Eficiencia en Cúbits / Recursos} & \textbf{Escalabilidad en la Era NISQ} \\
        \midrule
        \textbf{Gurobi Optimizer} \cite{Gurobi2025} & Clásico Determinista Exacto & Planos de corte y bifurcación ramificada. & Memoria RAM y CPU con escalado exponencial en peor caso. & Nula para instancias densas de gran escala ($N \gg 100$). \\
        \addlinespace
        \textbf{Simulated Annealing} & Clásico Estocástico Heurístico & Penalización o rechazo analítico de vecindarios. & Cómputo CPU secuencial ligero y altamente paralelizable. & Alta rapidez temporal, pero sin certificado de optimalidad global. \\
        \addlinespace
        \textbf{Standard QAOA} \cite{Farhi2014} & Cuántico Digital Variacional & Penalización cuadrática en la función de coste (QUBO). & Requiere $N$ cúbits sin registros auxiliares. & Baja (alta sensibilidad a mesetas estériles e inestabilidad por penalización). \\
        \addlinespace
        \textbf{Quantum Annealing} & Cuántico Adiabático Nativo & Mapeo de penalizaciones QUBO en grafos físicos. & Alta dependencia del grafo de acoplamiento del chip. & Media (condicionada por la conectividad limitada del hardware). \\
        \addlinespace
        \textbf{XY-QAOA Regularizado (Propuesto)} & Cuántico Híbrido Variacional & Confinamiento físico en subespacio combinatorio reducido. & Optimizado a $N$ cúbits lineales libres de registros auxiliares. & Alta (menor profundidad de circuito y mayor resistencia al ruido físico \cite{Barkoutsos2020}). \\
        \bottomrule
    \end{tabular}
\end{table}